% Chapter introduction: why the topic arose, who first asked, what it answers.
% STATUS: written 2026-11-04 (Bresar 1991 read in full)
A derivation follows the product rule. Several lines of work weaken that rule,
and each weakening starts from a different question. The first line began with
a question about norms. Brešar asked how far the composition of two
derivations sits from a wider class of maps, and he called that class the
generalized derivations \cite[p.~89]{Bresar1991Glasgow}. A second line keeps
the rule but turns it: multiplication on one side is twisted by an
endomorphism, and the twisted derivations arise. Each of these classes then
developed a life of its own, with its own Posner-type questions, its own
commutativity theorems, and its own operator examples.

This chapter follows the lines in turn. Section
\ref{sec:generalized-twisted:generalized} reports the generalized derivations,
from the norm estimate where they were named to the algebraic study of their
products and commutator identities. The sections after it take up the twists.
\hbfigure{06-generalized-twisted/specializations}{Twisted derivations and their specializations. With $\tau=\mathrm{id}$ one gets the $\sigma$-derivations of Ore; with $\sigma=\tau=\mathrm{id}$, ordinary derivations.}{fig:twisted:spec}
